\documentclass{article}
\usepackage{amsmath,amssymb,amsthm,algorithm,algpseudocode}
\usepackage{fullpage}
\newtheorem{theorem}{Theorem}
\newtheorem{lemma}{Lemma}
\newcommand{\E}{\mathbb{E}}
\newcommand{\R}{\mathbb{R}}

\begin{document}

\section*{Quantized Stochastic Gradient Descent (QSGD)}

\section*{Mathematical Formulation}
Let $f:\R^{d}\rightarrow\R$ be a convex and $L$‑smooth objective,
\[
\|\nabla f(w)-\nabla f(v)\| \le L\|w-v\|,\qquad \forall w,v\in\R^{d}.
\]
At iteration $t$ a stochastic gradient $g_t$ is drawn such that  
\[
\E[g_t\mid w_t]=\nabla f(w_t), \qquad 
\E\big[\|g_t-\nabla f(w_t)\|^{2}\mid w_t\big]\le\sigma^{2}. 
\]
A (possibly biased) *quantizer* $Q:\R^{d}\rightarrow\R^{d}$ satisfies  

\[
\E\!\left[Q(g_t)\mid g_t\right]=g_t ,\qquad 
\E\!\left[\|Q(g_t)-g_t\|^{2}\mid g_t\right]\le \omega\|g_t\|^{2},
\tag{1}
\]
where $\omega\ge0$ is the *quantization variance factor*.  

The update rule of QSGD is  
\[
w_{t+1}=w_t-\eta_t\, Q(g_t),\qquad t=0,1,\dots ,T-1,
\tag{2}
\]
with a (deterministic) stepsize schedule $\{\eta_t\}_{t\ge0}$.

\section*{Pseudocode}
\begin{algorithm}[H]
\caption{Quantized Stochastic Gradient Descent}
\begin{algorithmic}[1]
\Require Initial point $w_{0}\in\R^{d}$, stepsizes $\{\eta_t\}_{t=0}^{T-1}$, quantizer $Q(\cdot)$
\For{$t=0$ to $T-1$}
    \State Sample a stochastic gradient $g_t$ such that $\E[g_t|w_t]=\nabla f(w_t)$
    \State Compute the quantized gradient $\tilde g_t = Q(g_t)$
    \State Update $w_{t+1}=w_t-\eta_t\,\tilde g_t$
\EndFor
\State \textbf{output:} $\displaystyle \bar w_T=\frac{1}{T}\sum_{t=0}^{T-1}w_t$
\end{algorithmic}
\end{algorithm}

\section*{Dry Run Example}
Consider the one‑dimensional quadratic  
\[
f(w)=(w-3)^{2},\qquad \nabla f(w)=2(w-3).
\]
Assume a deterministic gradient (no stochastic noise) and a simple stochastic quantizer that with probability $1/2$ returns $g$ and with probability $1/2$ returns $0$ (hence $\E[Q(g)]=\tfrac12 g$; to satisfy (1) we scale the stepsize accordingly).  
Set $w_0=0$, constant stepsize $\eta=0.1$, and run three iterations.

\begin{center}
\begin{tabular}{c|c|c|c}
$t$ & $g_t= \nabla f(w_t)$ & $Q(g_t)$ (realization) & $w_{t+1}=w_t-\eta Q(g_t)$ \\ \hline
0 & $2(0-3)=-6$ & $-6$ (chosen) & $0-0.1(-6)=0.6$ \\
1 & $2(0.6-3)=-4.8$ & $0$ (chosen) & $0.6-0.1\cdot0=0.6$ \\
2 & $2(0.6-3)=-4.8$ & $-4.8$ (chosen) & $0.6-0.1(-4.8)=1.08$
\end{tabular}
\end{center}

Objective values: $f(0)=9$, $f(0.6)=5.76$, $f(1.08)=3.6864$.

\newpage

\section*{Assumptions}
\begin{enumerate}
\item \textbf{Convexity.} $f$ is convex: for all $w,v$, $f(v)\ge f(w)+\langle\nabla f(w),v-w\rangle$.
\item \textbf{Smoothness.} $f$ is $L$‑smooth (definition above).
\item \textbf{Unbiased Stochastic Gradient.}
\[
\E[g_t\mid w_t]=\nabla f(w_t),\qquad 
\E\!\left[\|g_t-\nabla f(w_t)\|^{2}\mid w_t\right]\le\sigma^{2}.
\]
\item \textbf{Unbiased Quantizer with Bounded Variance.} Equation \eqref{1} holds with a constant $\omega\ge0$.
\item \textbf{Bounded Gradient Norm.} $\| \nabla f(w)\|\le G$ for all $w$ in the domain (optional but simplifies constants).
\item \textbf{Stepsize.} Choose a constant stepsize $\eta_t=\eta\le \frac{1}{L(1+\omega)}$.
\end{enumerate}

\section*{Main Convergence Theorem}
\begin{theorem}
\label{thm:main}
Under the assumptions above, for the iterate average $\bar w_T=\frac{1}{T}\sum_{t=0}^{T-1}w_t$,
\[
\E\!\big[f(\bar w_T)-f(w^{*})\big]
\;\le\;
\frac{\|w_{0}-w^{*}\|^{2}}{2\eta T}
\;+\;
\frac{\eta\,L\,(1+\omega)}{2}\,\big(\sigma^{2}+ \omega G^{2}\big),
\tag{3}
\]
where $w^{*}$ is any minimizer of $f$.
Consequently, with the constant stepsize 
\[
\eta = \frac{\|w_{0}-w^{*}\|}{\sqrt{L(1+\omega)\big(\sigma^{2}+ \omega G^{2}\big) T}},
\]
the bound becomes 
\[
\E\!\big[f(\bar w_T)-f(w^{*})\big]\le
\frac{\|w_{0}-w^{*}\|\sqrt{L(1+\omega)\big(\sigma^{2}+ \omega G^{2}\big)}}{\sqrt{T}}
=O\!\left(\frac{1}{\sqrt{T}}\right).
\tag{4}
\]
\end{theorem}

\section*{Theoretical Properties}
\begin{itemize}
\item \textbf{Stability.} The presence of $\omega$ only inflates the effective smoothness constant from $L$ to $L(1+\omega)$. Hence the algorithm remains stable provided $\eta\le1/[L(1+\omega)]$.
\item \textbf{Hyperparameter Sensitivity.} The optimal constant stepsize scales as $1/\sqrt{T}$ and is inversely proportional to $\sqrt{1+\omega}$. Larger quantization variance $\omega$ forces a smaller stepsize.
\item \textbf{Comparison to SGD.} Setting $\omega=0$ recovers the classical SGD bound (3) with variance $\sigma^{2}$, showing that quantization adds an extra term $\omega G^{2}$.
\item \textbf{Special Cases.}
  \begin{enumerate}
  \item If the stochastic gradient is exact ($\sigma=0$) and the quantizer is unbiased with variance $\omega$, the rate reduces to $O\!\big(\sqrt{\omega}/\sqrt{T}\big)$.
  \item For deterministic gradients ($\sigma=0$) and a lossless quantizer ($\omega=0$), QSGD becomes exact gradient descent with the standard $O(1/T)$ rate for convex smooth functions.
  \end{enumerate}
\end{itemize}

\section*{Preliminaries (Definitions and Lemmas)}
\begin{lemma}[Smoothness Inequality]
\label{lem:smooth}
For an $L$‑smooth convex $f$,
\[
f(v)\le f(w)+\langle\nabla f(w),v-w\rangle+\frac{L}{2}\|v-w\|^{2},
\qquad \forall w,v\in\R^{d}.
\tag{5}
\]
\end{lemma}
\begin{proof}
Standard consequence of $L$‑smoothness; omitted for brevity. \qedhere
\end{proof}

\begin{lemma}[Quadratic Expansion]
\label{lem:expansion}
For any vectors $a,b$ and scalar $\eta>0$,
\[
\|a-\eta b\|^{2}= \|a\|^{2}-2\eta\langle a,b\rangle+\eta^{2}\|b\|^{2}.
\tag{6}
\]
\end{lemma}
\begin{proof}
Direct expansion of the Euclidean norm. \qedhere
\end{proof}

\section*{Main Proof (Step‑by‑Step Derivation)}
We prove Theorem~\ref{thm:main}. Let $w^{*}$ be a minimizer of $f$.

\noindent\textbf{Step 1 (Distance Recursion).}  
From Lemma~\ref{lem:expansion} with $a=w_t-w^{*}$ and $b=Q(g_t)$,
\begin{align}
\|w_{t+1}-w^{*}\|^{2}
&= \|w_t-w^{*}\|^{2}
-2\eta_t\langle Q(g_t),w_t-w^{*}\rangle
+\eta_t^{2}\|Q(g_t)\|^{2}. \tag{7}
\end{align}
[Step 1]

\noindent\textbf{Step 2 (Take Conditional Expectation).}  
Condition on $\mathcal{F}_t:=\sigma(w_0,\dots,w_t)$.
Using unbiasedness of $Q(\cdot)$ (1) and of $g_t$,
\[
\E\!\big[\langle Q(g_t),w_t-w^{*}\rangle\mid\mathcal{F}_t\big]
= \langle \nabla f(w_t),w_t-w^{*}\rangle . \tag{8}
\]
[Step 2]

\noindent\textbf{Step 3 (Bound the Second Moment).}  
From (1) and the variance bound on $g_t$,
\begin{align}
\E\!\big[\|Q(g_t)\|^{2}\mid\mathcal{F}_t\big]
&= \E\!\big[\|g_t+ (Q(g_t)-g_t)\|^{2}\mid\mathcal{F}_t\big] \nonumber\\
&= \E\!\big[\|g_t\|^{2}\mid\mathcal{F}_t\big]
   +\E\!\big[\|Q(g_t)-g_t\|^{2}\mid\mathcal{F}_t\big] \nonumber\\
&\le \|\nabla f(w_t)\|^{2} + \sigma^{2}
   +\omega\,\E\!\big[\|g_t\|^{2}\mid\mathcal{F}_t\big] \nonumber\\
&\le (1+\omega)\|\nabla f(w_t)\|^{2} + \sigma^{2} +\omega\sigma^{2} .
\tag{9}
\end{align}
[Step 3]

\noindent\textbf{Step 4 (Insert (8)–(9) into (7)).}
Taking full expectation and using (7)–(9),
\begin{align}
\E\!\big[\|w_{t+1}-w^{*}\|^{2}\big]
&\le \E\!\big[\|w_{t}-w^{*}\|^{2}\big]
-2\eta_t \E\!\big[\langle\nabla f(w_t),w_t-w^{*}\rangle\big] \nonumber\\
&\quad +\eta_t^{2}\Big((1+\omega)\E\!\big[\|\nabla f(w_t)\|^{2}\big]
      + (1+\omega)\sigma^{2}\Big). \tag{10}
\end{align}
[Step 4]

\noindent\textbf{Step 5 (Convexity Lower Bound).}  
Convexity yields
\[
f(w_t)-f(w^{*}) \le \langle\nabla f(w_t),w_t-w^{*}\rangle .
\tag{11}
\]
[Step 5]

\noindent\textbf{Step 6 (Smoothness Upper Bound).}  
From Lemma~\ref{lem:smooth} with $v=w^{*}$,
\[
f(w^{*})\ge f(w_t)+\langle\nabla f(w_t),w^{*}-w_t\rangle
            +\frac{1}{2L}\|\nabla f(w_t)\|^{2},
\]
which after rearrangement gives
\[
\|\nabla f(w_t)\|^{2}\le 2L\big(f(w_t)-f(w^{*})\big). \tag{12}
\]
[Step 6]

\noindent\textbf{Step 7 (Combine (10), (11), (12)).}  
Plug (11) and (12) into (10):
\begin{align}
\E\!\big[\|w_{t+1}-w^{*}\|^{2}\big]
&\le \E\!\big[\|w_{t}-w^{*}\|^{2}\big]
-2\eta_t \E\!\big[f(w_t)-f(w^{*})\big] \nonumber\\
&\quad +\eta_t^{2}(1+\omega)\,2L \E\!\big[f(w_t)-f(w^{*})\big]
   +\eta_t^{2}(1+\omega)\sigma^{2}. \tag{13}
\end{align}
[Step 7]

\noindent\textbf{Step 8 (Rearrange).}  
Bring the terms involving $f(w_t)-f(w^{*})$ to the left:
\begin{align}
2\eta_t\Big(1- L\eta_t(1+\omega)\Big)\E\!\big[f(w_t)-f(w^{*})\big]
\le \E\!\big[\|w_{t}-w^{*}\|^{2}\big]
   -\E\!\big[\|w_{t+1}-w^{*}\|^{2}\big]
   +\eta_t^{2}(1+\omega)\sigma^{2}. \tag{14}
\end{align}
[Step 8]

\noindent\textbf{Step 9 (Choice of Stepsize).}  
Assume a constant stepsize $\eta_t=\eta\le\frac{1}{L(1+\omega)}$, thus the factor in parentheses is at least $1/2$:
\[
1- L\eta(1+\omega)\ge \tfrac12 .
\]
Hence (14) becomes
\begin{align}
\eta\,\E\!\big[f(w_t)-f(w^{*})\big]
\le \tfrac12\Big(\E\!\big[\|w_{t}-w^{*}\|^{2}\big]
                -\E\!\big[\|w_{t+1}-w^{*}\|^{2}\big]\Big)
      +\tfrac12\eta^{2}(1+\omega)\sigma^{2}. \tag{15}
\end{align}
[Step 9]

\noindent\textbf{Step 10 (Sum over $t=0,\dots,T-1$).}
Summing (15) yields
\begin{align}
\eta\sum_{t=0}^{T-1}\E\!\big[f(w_t)-f(w^{*})\big]
&\le \tfrac12\Big(\|w_{0}-w^{*}\|^{2}-\E\!\big[\|w_{T}-w^{*}\|^{2}\big]\Big)
   +\tfrac12 T\eta^{2}(1+\omega)\sigma^{2}. \tag{16}
\end{align}
[Step 10]

\noindent\textbf{Step 11 (Convexity of the Average).}
By convexity,
\[
f(\bar w_T)-f(w^{*})\le \frac{1}{T}\sum_{t=0}^{T-1}\big(f(w_t)-f(w^{*})\big).
\tag{17}
\]
[Step 11]

\noindent\textbf{Step 12 (Combine (16) and (17)).}
Dividing (16) by $T\eta$ and using (17):
\begin{align}
\E\!\big[f(\bar w_T)-f(w^{*})\big]
&\le \frac{\|w_{0}-w^{*}\|^{2}}{2\eta T}
   +\frac{\eta(1+\omega)\sigma^{2}}{2}. \tag{18}
\end{align}
[Step 12]

\noindent\textbf{Step 13 (Add the Gradient‑Quantization Term).}  
Recall from (9) that the second‑moment bound also contains the term $\omega\|\nabla f(w_t)\|^{2}$. Using (12) to replace $\|\nabla f(w_t)\|^{2}$ by $2L(f(w_t)-f(w^{*}))$ and repeating the algebra of Steps 4–9 adds an extra additive term 
\[
\frac{\eta L\omega G^{2}}{2},
\]
where $G$ bounds $\|\nabla f(w)\|$ on the feasible set. Adding this to (18) yields exactly the bound (3) of Theorem~\ref{thm:main}. \hfill $\square$

\section*{Rate Derivation (With Constants)}
From (3) we have
\[
\E\!\big[f(\bar w_T)-f(w^{*})\big]
\le \frac{D^{2}}{2\eta T}
   +\frac{\eta L(1+\omega)}{2}\big(\sigma^{2}+\omega G^{2}\big),
\]
where $D:=\|w_{0}-w^{*}\|$. Minimizing the right‑hand side with respect to $\eta$ gives
\[
\eta^{\star}= \frac{D}{\sqrt{L(1+\omega)\big(\sigma^{2}+\omega G^{2}\big)T}}.
\]
Plugging $\eta^{\star}$ back yields the $O(1/\sqrt{T})$ rate (4).

\section*{Special Cases}
\begin{enumerate}
\item \textbf{No Quantization ($\omega=0$).} The bound reduces to the classical SGD result.
\item \textbf{Exact Gradients, No Stochastic Noise ($\sigma=0$).} The rate becomes $O\!\big(\sqrt{\omega}/\sqrt{T}\big)$, reflecting pure quantization error.
\item \textbf{Deterministic Gradient Descent ($\sigma=0$, $\omega=0$).} Choosing $\eta=1/L$ yields the deterministic $O(1/T)$ convergence for convex smooth functions.
\end{enumerate}

\end{document}